\begin{frame}{스트라이드 $s=2$가 실제로 하는 일}
  \begin{itemize}
    \item 1D에서 길이를 반으로, 2D에서는 \textbf{가로·세로 둘 다} 반으로
    \item $\Rightarrow$ 공간 \textbf{면적은 1/4}로 줄어든다
    \item 예: 224$\times$224 에 $k{=}3, p{=}1, s{=}2$ 이면
      $n_{\text{out}} = \lfloor(224+2-3)/2\rfloor+1 = \mathbf{112}$
  \end{itemize}
  \vskip0.6em
  \begin{block}{AlexNet(2012)의 첫 합성곱 층}
    정확히 이 설정(224$\to$112)으로 다운샘플링한다.
  \end{block}
  \vskip0.4em
  10.2절에서 2$\times$2 \textbf{풀링}도 면적을 1/4로 줄여주는 것을
  보게 되는데, ``둘 중 무엇을 쓰느냐''는 질문이 생긴다 ---
  이 절 FAQ에서 비교한다.
\end{frame}
